% ===============================================================
% Ficha 2 - Construção do DER
% Documento para preenchimento em grupo
% ===============================================================

\documentclass[11pt,a4paper]{article}

\newcommand{\coursecode}{CIC0101}
\newcommand{\coursename}{Sistemas de Informação}
\newcommand{\lecnumber}{2}
\newcommand{\lectitle}{Construção do DER}
\newcommand{\lecturer}{Patrícia Araújo de Oliveira}
\newcommand{\semester}{2026 / 2}
\newcommand{\groupnumber}{01}

% ------------- For Encoding & fonts ----------------
\usepackage[T1]{fontenc}
\usepackage[utf8]{inputenc}
\usepackage[brazil]{babel}
\usepackage{lmodern}
\usepackage[scaled=.92]{helvet}   % sans-serif for headings
\usepackage{microtype}

% ---------------- For Layout ----------------
\usepackage[margin=1in, headheight=15pt]{geometry}
\usepackage{setspace}
\setstretch{1.1}
\usepackage{enumitem}
\setlist{itemsep=2pt, topsep=4pt}
\usepackage{booktabs, array, tabularx, graphicx}
\usepackage{caption, subcaption}
\captionsetup{font=small, labelfont={bf,sf}, format=hang}

% ---------------- For Color ----------------
\usepackage[dvipsnames,svgnames]{xcolor}

% ---------------- For Math ----------------
\usepackage{amsmath, amssymb, amsthm, mathtools, bm}

% ---------------- For Boxes ----------------
\usepackage[most]{tcolorbox}

% ------------- For Code & algorithms ----------------
\usepackage{listings}
\usepackage{algorithm}
\usepackage[noend]{algpseudocode}

% -------------- For Figures / plots ----------------
\usepackage{tikz}
\usetikzlibrary{arrows.meta, positioning, shapes.geometric, calc, fit, backgrounds}
\usepackage{pgfplots}
\pgfplotsset{compat=newest}

% ------------- For Headings, headers ----------------
\usepackage{titlesec}
\usepackage{fancyhdr}
\usepackage{lastpage}

% -------- For Hyperlinks (load last) ----------------
\usepackage[colorlinks=true, linkcolor=accent, citecolor=accent3, urlcolor=accent2]{hyperref}

% ---------------- Colors (change the theme here) ----------------
\definecolor{accent}{HTML}{1F5FBF}     % main accent (blue)
\definecolor{accent2}{HTML}{C2410C}    % secondary accent (orange)
\definecolor{accent3}{HTML}{15803D}    % tertiary accent (green)
\definecolor{codebg}{HTML}{F6F8FA}
\definecolor{codegray}{HTML}{6A737D}
\definecolor{codekw}{HTML}{CF222E}
\definecolor{codestr}{HTML}{0A3069}
\definecolor{codefn}{HTML}{8250DF}

% ---------------- Section headings ----------------
\titleformat{\section}{\Large\sffamily\bfseries\color{accent}}{\thesection}{0.7em}{}
\titleformat{\subsection}{\large\sffamily\bfseries\color{black!80}}{\thesubsection}{0.7em}{}
\titleformat{\subsubsection}{\normalsize\sffamily\bfseries\color{black!70}}{\thesubsubsection}{0.7em}{}
\titlespacing*{\section}{0pt}{16pt}{8pt}

% ---------------- Header / footer ----------------
% \pagestyle{fancy}
% \fancyhf{}
% \fancyhead[L]{\sffamily\small\color{gray}\coursecode\ --- \coursename}
% \fancyhead[R]{\sffamily\small\color{gray}Ficha \lecnumber: \lectitle}
% \fancyfoot[C]{\sffamily\small\thepage}
% \renewcommand{\headrulewidth}{0.4pt}
% \renewcommand{\headrule}{\hbox to\headwidth{\color{gray!50}\leaders\hrule height \headrulewidth\hfill}}
\pagestyle{fancy}
\fancyhf{}
\fancyhead[L]{\sffamily\small\color{gray}\coursecode\ --- \coursename}
\fancyhead[R]{\sffamily\small\color{gray}Ficha \lecnumber: \lectitle}
\fancyfoot[L]{\sffamily\small\color{gray}Grupo \groupnumber}
\fancyfoot[C]{\sffamily\small\color{gray}\coursecode\ --- \coursename}
\fancyfoot[R]{\sffamily\small\color{gray}Página \thepage\ de \pageref{LastPage}}
\renewcommand{\headrulewidth}{0.4pt}
\renewcommand{\footrulewidth}{0.4pt}
\renewcommand{\headrule}{\hbox to\headwidth{\color{gray!50}\leaders\hrule height \headrulewidth\hfill}}
\renewcommand{\footrule}{\hbox to\headwidth{\color{gray!50}\leaders\hrule height \footrulewidth\hfill}}

% ---------------- Title block ----------------
\newcommand{\makelecturetitle}{%
  \begin{tcolorbox}[enhanced, colback=accent, colframe=accent, coltext=white,
    sharp corners, boxrule=0pt, left=14pt, right=14pt, top=12pt, bottom=12pt,
    overlay={\fill[white,opacity=0.08] (frame.south east) circle (2.2cm);}]
    {\sffamily\small Universidade de Bras\'ilia\hfill \semester\par}
    {\sffamily\small Departamento de Ci\^encia da Computa\c{c}\~ao\par}
    \vspace{4pt}
    {\sffamily\small\textbf{\coursecode}\ $\cdot$\ \coursename\par}
    \vspace{6pt}
    {\sffamily\bfseries\huge Ficha \lecnumber: \lectitle\par}
    \vspace{6pt}
    {\sffamily\small Professora: \lecturer\hfill Compilado: \today\par}
  \end{tcolorbox}
  \vspace{6pt}
}

% ---------------- Math macros ----------------
\newcommand{\R}{\mathbb{R}}
\newcommand{\N}{\mathbb{N}}
\newcommand{\Z}{\mathbb{Z}}
\newcommand{\E}{\mathbb{E}}
\renewcommand{\Pr}{\mathbb{P}}
% \O already exists in LaTeX (the letter O-slash), so keep it for text mode
% and use \mathcal{O} (Big-O) in math mode.
\let\oslashletter\O
\renewcommand{\O}{\ifmmode\mathcal{O}\else\oslashletter\fi}
\newcommand{\norm}[1]{\left\lVert #1 \right\rVert}
\newcommand{\abs}[1]{\left\lvert #1 \right\rvert}
\newcommand{\inner}[2]{\left\langle #1, #2 \right\rangle}
\DeclareMathOperator*{\argmin}{arg\,min}
\DeclareMathOperator*{\argmax}{arg\,max}

% Nicer QED symbol
\renewcommand{\qedsymbol}{$\blacksquare$}

% Inline code:  \code{merge\_sort(A)}
\newcommand{\code}[1]{\colorbox{codebg}{\texttt{\small #1}}}

\newcommand{\grupointegrantes}[2]{%
  \begin{tcolorbox}[colback=white!92!accent, colbacktitle=accent,
    colframe=accent, coltitle=white, coltext=black,
    title={\sffamily\bfseries #1 --- Integrantes},
    boxrule=0.6pt, arc=1pt, left=14pt, right=14pt, top=6pt, bottom=6pt,
    toptitle=4pt, bottomtitle=4pt]
    \small
    \begin{tabularx}{\linewidth}{@{}>{\raggedright\arraybackslash}X r
                                    >{\raggedright\arraybackslash}X r@{}}
      #2
    \end{tabularx}
  \end{tcolorbox}%
  \vspace{8pt}%
}

% ---------------- Base style for all boxes ----------------
\tcbset{
  thmbase/.style={
    enhanced, breakable, sharp corners, boxrule=0pt, frame hidden,
    left=10pt, right=8pt, top=6pt, bottom=6pt,
    fonttitle=\sffamily\bfseries, before skip=10pt, after skip=10pt
  }
}

% ---------------- Numbered boxes ----------------
% Usage: \begin{theorem}{Title}{label} ... \end{theorem}   ref: Theorem~\ref{thm:label}
\newtcbtheorem[number within=section]{theorem}{Theorem}%
  {thmbase, colback=accent!6, borderline west={3pt}{0pt}{accent}, coltitle=accent}{thm}
\newtcbtheorem[use counter from=theorem]{lemma}{Lemma}%
  {thmbase, colback=accent!6, borderline west={3pt}{0pt}{accent}, coltitle=accent}{lem}
\newtcbtheorem[use counter from=theorem]{definition}{Definition}%
  {thmbase, colback=accent3!7, borderline west={3pt}{0pt}{accent3}, coltitle=accent3}{def}
\newtcbtheorem[use counter from=theorem]{example}{Example}%
  {thmbase, colback=accent2!6, borderline west={3pt}{0pt}{accent2}, coltitle=accent2}{ex}
\newtcbtheorem[number within=section]{exercise}{Exercise}%
  {thmbase, colback=gray!8, borderline west={3pt}{0pt}{gray}, coltitle=black!70}{exc}

% ---------------- Unnumbered boxes ----------------
\newtcolorbox{keyidea}[1][Key Idea]{thmbase, colback=yellow!12,
  borderline west={3pt}{0pt}{yellow!70!orange},
  before upper={\textbf{\sffamily\color{yellow!40!black}#1.}\ }}
\newtcolorbox{notebox}[1][Note]{thmbase, colback=gray!8,
  borderline west={3pt}{0pt}{gray!70},
  before upper={\textbf{\sffamily\color{black!70}#1.}\ }}
\newtcolorbox{warnbox}[1][Warning]{thmbase, colback=red!6,
  borderline west={3pt}{0pt}{red!70!black},
  before upper={\textbf{\sffamily\color{red!70!black}#1.}\ }}
\newtcolorbox{objectives}{thmbase, colback=accent!6,
  borderline west={3pt}{0pt}{accent},
  before upper={\textbf{\sffamily\color{accent}Learning objectives}\par\vspace{2pt}}}

% ---------------- Listings style ----------------
\lstdefinestyle{modern}{
  backgroundcolor=\color{codebg},
  basicstyle=\ttfamily\small,
  commentstyle=\itshape\color{codegray},
  keywordstyle=\bfseries\color{codekw},
  stringstyle=\color{codestr},
  identifierstyle=\color{black},
  numbers=left, numberstyle=\tiny\color{codegray}, numbersep=8pt,
  frame=lines, rulecolor=\color{gray!40},
  framexleftmargin=4pt, xleftmargin=14pt,
  breaklines=true, breakatwhitespace=true,
  postbreak=\mbox{\textcolor{codegray}{$\hookrightarrow$}\space},
  showstringspaces=false, tabsize=4, captionpos=b,
  aboveskip=10pt, belowskip=6pt
}
\lstset{style=modern}

% ---------------- Algorithm style ----------------
\algrenewcommand\algorithmicrequire{\textbf{Input:}}
\algrenewcommand\algorithmicensure{\textbf{Output:}}
\floatname{algorithm}{Algorithm}


\begin{document}

\makelecturetitle
\vspace{2pt}

\grupointegrantes{Grupo \groupnumber}{
  \textcolor{accent}{Nome do integrante 1} & \textcolor{accent}{Matrícula} &
  \textcolor{accent}{Nome do integrante 2} & \textcolor{accent}{Matrícula} \\
  \textcolor{accent}{Nome do integrante 3} & \textcolor{accent}{Matrícula} &
  \textcolor{accent}{Nome do integrante 4} & \textcolor{accent}{Matrícula}
}

\begin{notebox}[Ficha 2]
Aula de 29/10/2026 \,$\cdot$\, Projeto Integrador \,$\cdot$\, Trabalho em grupo.
\end{notebox}

\begin{objectives}
Representar graficamente o MER elaborado na aula anterior, utilizando uma notação consistente.
\end{objectives}

\tableofcontents
\vspace{8pt}

\section{Consolidar entidades}

Para cada entidade, descrevam brevemente o que ela representa no sistema e verifiquem se
atributos e identificadores estão definidos. Consultem a atividade assíncrona, o BPMN e os
requisitos funcionais para identificar possíveis ajustes.

\begin{table}[!ht]
\centering
\small
\setlength{\tabcolsep}{5pt}
\renewcommand{\arraystretch}{1.8}
\begin{tabularx}{\linewidth}{|>{\raggedright\arraybackslash}p{4.5cm}
                              |>{\raggedright\arraybackslash}X|}
\hline
\textbf{Entidade} & \textbf{Papel} \\
\hline
\textbf{SOLICITAÇÃO (exemplo)} &
Registra o pedido de aproveitamento de estudos realizado pelo estudante, incluindo a data de
abertura e a situação da solicitação. \\
\hline
 & \\ \hline
 & \\ \hline
 & \\ \hline
 & \\ \hline
 & \\ \hline
 & \\ \hline
 & \\ \hline
 & \\ \hline
\end{tabularx}
\end{table}

\begin{warnbox}[Observação]
Se faltar linha, continue registrando as entidades na parte de trás da folha ou em uma folha adicional.
\end{warnbox}

\section{Cardinalidades e regras}

Para cada relacionamento, respondam: uma ocorrência de A se associa a quantas de B?
E o inverso? Justifiquem com BPMN, requisitos e regras do processo. Se não souberem,
registem a hipótese a validar.

\begin{table}[!ht]
\centering
\scriptsize
\setlength{\tabcolsep}{3pt}
\renewcommand{\arraystretch}{1.65}
\begin{tabularx}{\linewidth}{|>{\raggedright\arraybackslash}p{3.5cm}
                              |>{\centering\arraybackslash}X
                              |>{\centering\arraybackslash}X
                              |>{\raggedright\arraybackslash}X|}
\hline
\textbf{Relacionamento} &
\textbf{A $\rightarrow$ B} &
\textbf{B $\rightarrow$ A} &
\textbf{Regra de negócio / fonte} \\
\hline
\textbf{Estudante — realiza — Solicitação} &
1:N (hipótese) &
N:1 (hipótese) &
BPMN: atividade do estudante (enviar solicitação) \\
\hline
 & & & \\ \hline
 & & & \\ \hline
 & & & \\ \hline
 & & & \\ \hline
 & & & \\ \hline
 & & & \\ \hline
 & & & \\ \hline
 & & & \\ \hline
 & & & \\ \hline
 & & & \\ \hline
\end{tabularx}
\end{table}

\begin{warnbox}[Observação]
Se faltar linha, continue registrando os relacionamentos em uma folha adicional.
\end{warnbox}

\section{Espaço para o diagrama}

Desenhem o DER neste espaço ou anexem uma folha A4 com o desenho. Incluam entidades,
atributos, identificadores, relacionamentos e cardinalidades.

\begin{tcolorbox}[enhanced, colback=white, colframe=accent, boxrule=0.8pt,
  sharp corners, left=6pt, right=6pt, top=6pt, bottom=6pt, height=13cm]
\centering
\vfill
\color{gray}
\large Espaço reservado para o diagrama DER
\vfill
\end{tcolorbox}

\section{Verificação antes de concluir}

\begin{tcolorbox}[colback=accent!5, colframe=accent, boxrule=0.6pt, sharp corners]
\begin{tabularx}{\linewidth}{X X}
$\square$ Entidades e atributos do MER representados &
$\square$ Identificadores destacados \\
$\square$ Relacionamentos nomeados &
$\square$ Cardinalidades justificadas \\
$\square$ Todos os requisitos funcionais considerados &
$\square$ Dúvidas registradas \\
\end{tabularx}
\end{tcolorbox}

\subsection{Dúvidas ou regras ainda não confirmadas}
\begin{tcolorbox}[colback=white, colframe=gray!70, boxrule=0.6pt, sharp corners, height=9cm]
\end{tcolorbox}

\begin{notebox}[Entrega da oficina]
Ficha 1 (MER) + Ficha 2 (DER). Não é necessário implementar banco de dados.
\end{notebox}


\end{document}
